%============================================================
% rapport.tex
% Rapport explicatif : Tri Bitonique & Minage Blockchain en Mercury
%============================================================
\documentclass[12pt,a4paper]{article}

% ── Packages ────────────────────────────────────────────────
\usepackage[utf8]{inputenc}
\usepackage[T1]{fontenc}
\usepackage[french]{babel}
\usepackage{geometry}
\usepackage{graphicx}
\usepackage{amsmath,amssymb}
\usepackage{listings}
\usepackage{xcolor}
\usepackage{hyperref}
\usepackage{fancyhdr}
\usepackage{titlesec}
\usepackage{booktabs}
\usepackage{array}
\usepackage{tikz}
\usepackage{pgfplots}
\usepackage{algorithm}
\usepackage{algpseudocode}
\usepackage{mdframed}
\usepackage{tcolorbox}
\usepackage{multicol}
\usepackage{enumitem}
\usepackage{fontawesome5}

\pgfplotsset{compat=1.18}
\usetikzlibrary{shapes,arrows,positioning,calc,shadows,decorations.pathmorphing}

% ── Géométrie ───────────────────────────────────────────────
\geometry{
  top=2.5cm, bottom=2.5cm,
  left=2.5cm, right=2.5cm,
  headheight=15pt
}

% ── Couleurs ────────────────────────────────────────────────
\definecolor{mercuryblue}{RGB}{13,71,161}
\definecolor{mercurycyan}{RGB}{0,131,143}
\definecolor{blockgreen}{RGB}{27,94,32}
\definecolor{blockorange}{RGB}{230,81,0}
\definecolor{codebg}{RGB}{245,245,245}
\definecolor{codeframe}{RGB}{13,71,161}
\definecolor{keywordcolor}{RGB}{13,71,161}
\definecolor{commentcolor}{RGB}{56,142,60}
\definecolor{stringcolor}{RGB}{183,28,28}
\definecolor{lightblue}{RGB}{227,242,253}
\definecolor{lightgreen}{RGB}{232,245,233}
\definecolor{lightorange}{RGB}{255,243,224}

% ── Style listings (Mercury) ────────────────────────────────
\lstdefinelanguage{Mercury}{
  morekeywords={module,interface,implementation,import_module,
                pred,func,type,inst,mode,det,semidet,nondet,
                multi,erroneous,failure,is,if,then,else,true,
                fail,not,some,all,promise_pure,promise_semipure,
                pragma,foreign_code,foreign_proc,foreign_decl,
                foreign_import_module,initialise,finalise,
                mutable,use_module,include_module,end_module,
                where,typeclass,instance,lambda,fun,
                in,out,di,uo,mdi,muo,oi,bi,
                int,float,char,string,bool,list,io,
                yes,no,det,semidet,cc_nondet,cc_multi},
  sensitive=true,
  morecomment=[l]{\%},
  morestring=[b]",
  morestring=[b]'
}

\lstset{
  language=Mercury,
  backgroundcolor=\color{codebg},
  basicstyle=\ttfamily\footnotesize,
  keywordstyle=\color{keywordcolor}\bfseries,
  commentstyle=\color{commentcolor}\itshape,
  stringstyle=\color{stringcolor},
  numberstyle=\tiny\color{gray},
  numbers=left,
  numbersep=8pt,
  frame=single,
  framerule=1.2pt,
  rulecolor=\color{codeframe},
  breaklines=true,
  breakatwhitespace=false,
  tabsize=4,
  showstringspaces=false,
  captionpos=b,
  xleftmargin=15pt,
  xrightmargin=5pt
}

% ── En-tête / pied de page ──────────────────────────────────
\pagestyle{fancy}
\fancyhf{}
\fancyhead[L]{\small\color{mercuryblue}\textbf{Mercury — Tri Bitonique \& Minage Blockchain}}
\fancyhead[R]{\small\color{mercuryblue}INF322}
\fancyfoot[C]{\small\thepage}
\fancyfoot[L]{\small\color{gray}Université de Yaoundé I}
\fancyfoot[R]{\small\color{gray}\today}
\renewcommand{\headrulewidth}{1.5pt}
\renewcommand{\headrule}{\hbox to\headwidth{\color{mercuryblue}\leaders\hrule height \headrulewidth\hfill}}

% ── Titres ──────────────────────────────────────────────────
\titleformat{\section}
  {\Large\bfseries\color{mercuryblue}}
  {\thesection}{1em}{}
  [\color{mercuryblue}\titlerule]

\titleformat{\subsection}
  {\large\bfseries\color{mercurycyan}}
  {\thesubsection}{1em}{}

\titleformat{\subsubsection}
  {\normalsize\bfseries\color{blockorange}}
  {\thesubsubsection}{1em}{}

% ── Boîtes colorées ─────────────────────────────────────────
\tcbuselibrary{skins,breakable}

\newtcolorbox{definitionbox}[1]{
  colback=lightblue, colframe=mercuryblue,
  fonttitle=\bfseries\color{white},
  title={\faBook\quad Définition : #1},
  coltitle=white, attach boxed title to top left,
  boxed title style={colback=mercuryblue, rounded corners},
  rounded corners, breakable
}

\newtcolorbox{algorithmbox}[1]{
  colback=lightgreen, colframe=blockgreen,
  fonttitle=\bfseries\color{white},
  title={\faCogs\quad Algorithme : #1},
  coltitle=white, attach boxed title to top left,
  boxed title style={colback=blockgreen, rounded corners},
  rounded corners, breakable
}

\newtcolorbox{warningbox}{
  colback=lightorange, colframe=blockorange,
  fonttitle=\bfseries\color{white},
  title={\faExclamationTriangle\quad Remarque importante},
  coltitle=white, attach boxed title to top left,
  boxed title style={colback=blockorange, rounded corners},
  rounded corners
}

\newtcolorbox{resultbox}{
  colback=lightgreen, colframe=blockgreen,
  fonttitle=\bfseries\color{white},
  title={\faCheckCircle\quad Résultat},
  coltitle=white, rounded corners
}

% ── Commandes utiles ────────────────────────────────────────
\newcommand{\mercury}{\textsc{Mercury}}
\newcommand{\Oh}[1]{\mathcal{O}\!\left(#1\right)}
\newcommand{\code}[1]{\texttt{\color{keywordcolor}#1}}

%============================================================
\begin{document}

% ── PAGE DE TITRE ────────────────────────────────────────────
\begin{titlepage}
\begin{tikzpicture}[remember picture, overlay]
  % Fond dégradé simulé
  \fill[mercuryblue] (current page.north west) rectangle
        ([yshift=-4cm]current page.north east);
  \fill[mercurycyan!30] (current page.south west) rectangle
        ([yshift=2cm]current page.south east);
  % Cercles décoratifs
  \fill[white,opacity=0.05] ([xshift=12cm,yshift=-2cm]current page.north west)
        circle (5cm);
  \fill[white,opacity=0.05] ([xshift=2cm,yshift=-1cm]current page.north west)
        circle (3cm);
\end{tikzpicture}

\vspace*{2cm}
\begin{center}
  {\color{white}\Large\bfseries Université de Yaoundé I}\\[0.3cm]
  {\color{white}\large Département d'Informatique}\\[0.3cm]
  {\color{white}\normalsize INF322 — Programmation Logique et Fonctionnelle}
\end{center}

\vspace{2cm}
\begin{center}
  \begin{tcolorbox}[
    colback=white, colframe=white,
    width=0.85\textwidth,
    boxrule=0pt, arc=10pt,
    shadow={3mm}{-3mm}{0mm}{black!30}
  ]
    \begin{center}
      \vspace{0.5cm}
      {\huge\bfseries\color{mercuryblue}
        Tri Bitonique \&\\[0.3cm]
        Minage de Blockchain}\\[0.5cm]
      {\large\color{mercurycyan}
        Implémentation Experte en \mercury}\\[0.3cm]
      \hrule\vspace{0.3cm}
      {\normalsize\color{gray}
        Rapport Technique — Algorithmes Avancés en Programmation Logique}
      \vspace{0.5cm}
    \end{center}
  \end{tcolorbox}
\end{center}

\vspace{2cm}
\begin{center}
  \begin{tabular}{ll}
    \color{mercuryblue}\faCalendar\quad\textbf{Date} & \today \\[0.3cm]
    \color{mercuryblue}\faCode\quad\textbf{Langage} & Mercury 14.01+ \\[0.3cm]
    \color{mercuryblue}\faUniversity\quad\textbf{Cours} & INF322 \\[0.3cm]
    \color{mercuryblue}\faFlag\quad\textbf{Niveau} & L2 Informatique \\
  \end{tabular}
\end{center}
\vfill
\begin{center}
  \color{gray}\small Version 1.0 — Document généré automatiquement
\end{center}
\end{titlepage}

% ── TABLE DES MATIÈRES ───────────────────────────────────────
\tableofcontents
\newpage

%============================================================
\section{Introduction au langage Mercury}
%============================================================

\mercury{} est un langage de programmation logique-fonctionnel pur, développé
à l'Université de Melbourne (Australie) à partir de 1993 par Fergus Henderson
et Thomas Conway. Il combine la puissance expressive de Prolog avec la rigueur
des types de Haskell et les performances du C.

\subsection{Caractéristiques fondamentales}

\begin{multicols}{2}
\begin{itemize}[leftmargin=*]
  \item \textbf{Typage statique fort} : vérification à la compilation
  \item \textbf{Modes et déterminisme} : \code{det}, \code{semidet}, \code{nondet}
  \item \textbf{Pureté} : pas d'effets de bord cachés
  \item \textbf{Unification} : mécanisme fondamental de résolution
  \item \textbf{Backtracking} : retour arrière automatique
  \item \textbf{Performance} : compilation native C/Java/C\#
  \item \textbf{FFI} : interface avec C pour les bibliothèques natives
  \item \textbf{Modules} : système de modules complet
\end{itemize}
\end{multicols}

\subsection{Structure d'un programme Mercury}

\begin{lstlisting}[caption={Structure de base d'un module Mercury}]
:- module nom_module.
:- interface.
:- import_module io.

% Declarations publiques
:- pred predicat(type::mode) is determinisme.

:- implementation.
:- import_module list, int, string.

% Definitions des predicats
predicat(Param) :-
    corps_du_predicat(Param).
\end{lstlisting}

\begin{warningbox}
En \mercury{}, chaque prédicat doit avoir un \textbf{mode} déclaré
(\code{in}/\code{out}/\code{di}/\code{uo}) et un \textbf{déterminisme}
(\code{det}, \code{semidet}, \code{nondet}, \code{multi}).
L'état d'entrée-sortie \code{io::di, io::uo} garantit la pureté des effets.
\end{warningbox}

%============================================================
\section{Partie I : Tri Bitonique}
%============================================================

\subsection{Définition et propriétés}

\begin{definitionbox}{Séquence bitonique}
Une séquence $a_1, a_2, \ldots, a_n$ est dite \textbf{bitonique} si et seulement
si il existe un indice $k$ tel que :
$$a_1 \leq a_2 \leq \cdots \leq a_k \geq a_{k+1} \geq \cdots \geq a_n$$
ou si la séquence peut être obtenue par rotation cyclique d'une telle séquence.
En d'autres termes : la séquence croît puis décroît (ou l'inverse).
\end{definitionbox}

\begin{definitionbox}{Tri bitonique}
Le \textbf{tri bitonique} (Bitonic Sort, Batcher 1968) est un algorithme de tri
basé sur des réseaux de comparateurs. Il trie une séquence en :
\begin{enumerate}
  \item Construisant récursivement une séquence bitonique par fusion de
        sous-séquences croissantes et décroissantes alternées.
  \item Fusionnant la séquence bitonique obtenue dans l'ordre désiré.
\end{enumerate}
\end{definitionbox}

\subsection{Complexité}

\begin{center}
\begin{tabular}{lcccc}
\toprule
\textbf{Algorithme} & \textbf{Comparaisons} & \textbf{Étapes parallèles}
                    & \textbf{Stable} & \textbf{En place} \\
\midrule
Tri bitonique    & $\Oh{n \log^2 n}$ & $\Oh{\log^2 n}$ & Non & Oui \\
Tri fusion       & $\Oh{n \log n}$   & $\Oh{n}$        & Oui & Non \\
Tri rapide       & $\Oh{n \log n}$   & $\Oh{n}$        & Non & Oui \\
Tri bulles       & $\Oh{n^2}$        & $\Oh{n}$        & Oui & Oui \\
\bottomrule
\end{tabular}
\end{center}

\begin{warningbox}
Le tri bitonique n'est \textbf{pas} le plus efficace séquentiellement
$\Oh{n\log^2 n}$ vs $\Oh{n\log n}$), mais il est \textbf{optimal pour la
parallélisation} : ses $\Oh{\log^2 n}$ étapes parallèles en font le choix
privilégié pour les GPU, FPGA et processeurs vectoriels.
\end{warningbox}

\subsection{Principe de fonctionnement}

\begin{algorithmbox}{Tri bitonique récursif}
\begin{algorithm}[H]
\caption{Tri bitonique}
\begin{algorithmic}[1]
\Procedure{BitonicSort}{$A$, $dir$}
  \If{$|A| \leq 1$} \Return $A$ \EndIf
  \State $n \gets |A|$, $\text{mid} \gets n/2$
  \State $L \gets A[0..\text{mid}-1]$,\quad $R \gets A[\text{mid}..n-1]$
  \State $L' \gets $ \Call{BitonicSort}{$L$, \textsc{Croissant}}
  \State $R' \gets $ \Call{BitonicSort}{$R$, \textsc{Décroissant}}
  \State \Return \Call{BitonicMerge}{$L' \| R'$, $dir$}
\EndProcedure

\Procedure{BitonicMerge}{$A$, $dir$}
  \If{$|A| \leq 1$} \Return $A$ \EndIf
  \State $\text{mid} \gets |A|/2$
  \For{$i \gets 0$ \textbf{to} $\text{mid}-1$}
    \If{$dir = \textsc{Croissant}$ \textbf{and} $A[i] > A[i+\text{mid}]$}
      \State \Call{Échanger}{$A[i]$, $A[i+\text{mid}]$}
    \ElsIf{$dir = \textsc{Décroissant}$ \textbf{and} $A[i] < A[i+\text{mid}]$}
      \State \Call{Échanger}{$A[i]$, $A[i+\text{mid}]$}
    \EndIf
  \EndFor
  \State \Return \Call{BitonicMerge}{$A[0..\text{mid}-1]$, $dir$}
         $\|$ \Call{BitonicMerge}{$A[\text{mid}..n-1]$, $dir$}
\EndProcedure
\end{algorithmic}
\end{algorithm}
\end{algorithmbox}

\subsection{Illustration visuelle — Réseau de tri pour n=8}

\begin{center}
\begin{tikzpicture}[scale=0.85]
  % Fils
  \foreach \y in {0,1,2,3,4,5,6,7} {
    \draw[gray!50, thick] (0, \y) -- (14, \y);
    \node[left] at (0,\y) {\small $a_{\the\numexpr8-\y\relax}$};
  }

  % Étape 1 : construction bitonique - paires
  \foreach \y/\yy in {6/7,4/5,2/3,0/1} {
    \draw[mercuryblue,very thick] (1.5,\y) -- (1.5,\yy);
    \filldraw[mercuryblue] (1.5,\y) circle (3pt);
    \filldraw[mercuryblue] (1.5,\yy) circle (3pt);
  }

  % Étape 2
  \foreach \y/\yy in {4/7,5/6,0/3,1/2} {
    \draw[mercurycyan,very thick] (3,\y) -- (3,\yy);
    \filldraw[mercurycyan] (3,\y) circle (3pt);
    \filldraw[mercurycyan] (3,\yy) circle (3pt);
  }
  \foreach \y/\yy in {4/6,5/7,0/2,1/3} {
    \draw[mercurycyan!70,very thick] (4,\y) -- (4,\yy);
    \filldraw[mercurycyan!70] (4,\y) circle (3pt);
    \filldraw[mercurycyan!70] (4,\yy) circle (3pt);
  }

  % Étape 3 - fusion globale
  \foreach \y/\yy in {0/7,1/6,2/5,3/4} {
    \draw[blockgreen,very thick] (6,\y) -- (6,\yy);
    \filldraw[blockgreen] (6,\y) circle (3pt);
    \filldraw[blockgreen] (6,\yy) circle (3pt);
  }
  \foreach \y/\yy in {0/3,1/2,4/7,5/6} {
    \draw[blockgreen!80,very thick] (8,\y) -- (8,\yy);
    \filldraw[blockgreen!80] (8,\y) circle (3pt);
    \filldraw[blockgreen!80] (8,\yy) circle (3pt);
  }
  \foreach \y/\yy in {0/1,2/3,4/5,6/7} {
    \draw[blockgreen!60,very thick] (10,\y) -- (10,\yy);
    \filldraw[blockgreen!60] (10,\y) circle (3pt);
    \filldraw[blockgreen!60] (10,\yy) circle (3pt);
  }

  % Légende étapes
  \node[mercuryblue,above] at (1.5,7.5) {\tiny Étape 1};
  \node[mercurycyan,above] at (3.5,7.5) {\tiny Étape 2};
  \node[blockgreen,above] at (8,7.5) {\tiny Fusion bitonique};

  % Résultat
  \foreach \y in {0,1,2,3,4,5,6,7} {
    \node[right] at (14,\y) {\small $s_{\the\numexpr\y+1\relax}$};
  }
\end{tikzpicture}
\end{center}

\subsection{Implémentation Mercury — Prédicats clés}

\subsubsection{Prédicat principal \code{bitonic\_sort/3}}

\begin{lstlisting}[caption={Tri bitonique récursif en Mercury}]
:- pred bitonic_sort(direction::in, list(int)::in, list(int)::out) is det.

bitonic_sort(_Dir, [], []).
bitonic_sort(_Dir, [X], [X]).
bitonic_sort(Dir, List, Sorted) :-
    length(List, N),
    Half = N / 2,
    split(List, Half, Left, Right),
    % Trier la moitie gauche en croissant
    bitonic_sort(ascending,  Left,  SortedLeft),
    % Trier la moitie droite en decroissant
    bitonic_sort(descending, Right, SortedRight),
    % Concatener pour former la sequence bitonique
    append(SortedLeft, SortedRight, Bitonic),
    % Fusionner dans la direction souhaitee
    bitonic_merge(Dir, Bitonic, Sorted).
\end{lstlisting}

\subsubsection{Fusion bitonique \code{bitonic\_merge/3}}

\begin{lstlisting}[caption={Fusion bitonique par compare-and-swap}]
:- pred bitonic_merge(direction::in, list(int)::in, list(int)::out) is det.

bitonic_merge(Dir, List, Merged) :-
    length(List, N),
    ( N =< 1 -> Merged = List ;
        Half = N / 2,
        split(List, Half, Left, Right),
        compare_and_swap(Dir, Left, Right, NewLeft, NewRight),
        bitonic_merge(Dir, NewLeft,  MergedLeft),
        bitonic_merge(Dir, NewRight, MergedRight),
        append(MergedLeft, MergedRight, Merged)
    ).
\end{lstlisting}

\subsection{Trace d'exécution}

\begin{resultbox}
\begin{verbatim}
Entree  : [3, 7, 4, 8, 6, 2, 1, 5]

Phase 1 - Construction bitonique :
  [3,7] croissant  -> [3,7]
  [4,8] decroissant -> [8,4]
  -> Bitonic [3,7,8,4]  -> merge -> [3,4,7,8]

  [6,2] croissant  -> [2,6]
  [1,5] decroissant -> [5,1]
  -> Bitonic [2,6,5,1]  -> merge -> [1,2,5,6]

Phase 2 - Fusion finale :
  Bitonic [1,2,5,6,8,7,4,3] -> merge croissant

Sortie  : [1, 2, 3, 4, 5, 6, 7, 8]
Verification : OK
\end{verbatim}
\end{resultbox}

%============================================================
\section{Partie II : Minage de Blockchain}
%============================================================

\subsection{Fondements de la Blockchain}

\begin{definitionbox}{Blockchain}
Une \textbf{blockchain} (chaîne de blocs) est une structure de données
distribuée et immuable dans laquelle chaque bloc contient :
\begin{itemize}
  \item Un \textbf{identifiant} (index)
  \item Des \textbf{données} (transactions)
  \item Un \textbf{nonce} (Number used ONCE) pour le minage
  \item Un \textbf{hash} cryptographique du bloc courant
  \item Le \textbf{hash} du bloc précédent (\emph{prev\_hash})
\end{itemize}
Le chaînage par les hashs garantit l'intégrité : toute modification d'un
bloc invalide tous les blocs suivants.
\end{definitionbox}

\subsection{Preuve de Travail (Proof-of-Work)}

\begin{definitionbox}{Proof-of-Work (PoW)}
Le \textbf{Proof-of-Work} est le mécanisme de consensus qui rend le minage
coûteux en calcul. Le mineur doit trouver un nonce $n$ tel que :
$$\text{Hash}(\text{index} \| \text{données} \| n \| \text{prev\_hash})
  \prec \text{Cible}$$
où $\prec$ signifie que le hash commence par au moins $d$ zéros
(\emph{difficulté} $d$).
\end{definitionbox}

\subsection{Structure de données en Mercury}

\begin{lstlisting}[caption={Type bloc en Mercury}]
:- type block
    --->    block(
                b_index     :: int,      % Position dans la chaine
                b_data      :: string,   % Transactions
                b_nonce     :: int,      % Valeur trouvee par minage
                b_hash      :: string,   % Hash de ce bloc
                b_prev_hash :: string    % Hash du bloc precedent
            ).

:- type blockchain == list(block).
:- type difficulty == int.
\end{lstlisting}

\subsection{Algorithme de minage}

\begin{algorithmbox}{Proof-of-Work Mining}
\begin{algorithm}[H]
\caption{Minage d'un bloc}
\begin{algorithmic}[1]
\Procedure{MineBloc}{$\text{index}$, $\text{données}$, $\text{prev\_hash}$, $d$}
  \State $\text{nonce} \gets 0$
  \Repeat
    \State $s \gets \text{concat}(\text{index}, \text{données},
           \text{nonce}, \text{prev\_hash})$
    \State $h \gets \text{Hash}(s)$
    \If{$h$ commence par $d$ zéros}
      \Return $(\text{nonce}, h)$
    \EndIf
    \State $\text{nonce} \gets \text{nonce} + 1$
  \Until{solution trouvée}
\EndProcedure
\end{algorithmic}
\end{algorithm}
\end{algorithmbox}

\subsection{Implémentation Mercury — Minage}

\begin{lstlisting}[caption={Minage par recherche de nonce en Mercury}]
:- pred mine_nonce(int::in, string::in, string::in, difficulty::in,
                   int::in, int::out, string::out) is det.

mine_nonce(Index, Data, PrevHash, Diff, Try, Nonce, Hash) :-
    % Calculer le hash pour le nonce courant
    H = hash_block(Index, Data, Try, PrevHash),
    ( meets_difficulty(H, Diff) ->
        % Nonce trouve : solution valide
        Nonce = Try,
        Hash  = H
    ;
        % Essayer le nonce suivant (recursivite terminale)
        mine_nonce(Index, Data, PrevHash, Diff, Try + 1, Nonce, Hash)
    ).

% Verifier si le hash satisfait la difficulte
:- pred meets_difficulty(string::in, difficulty::in) is semidet.

meets_difficulty(Hash, Diff) :-
    Diff > 0,
    string.left(Hash, Diff) = Prefix,
    all_zeros(Prefix).
\end{lstlisting}

\subsection{Validation de l'intégrité de la chaîne}

\begin{lstlisting}[caption={Validation de la blockchain en Mercury}]
:- pred validate_chain(blockchain::in, bool::out) is det.

validate_chain([], yes).
validate_chain([_], yes).
validate_chain([B1, B2 | Rest], Valid) :-
    % Recalculer le hash du bloc B2
    ExpectedHash = hash_block(B2^b_index, B2^b_data,
                              B2^b_nonce, B2^b_prev_hash),
    % Verifier coherence : hash calcule = hash stocke
    %                  ET : hash de B2 = prev_hash de B1
    ( ExpectedHash = B2^b_hash,
      B2^b_hash = B1^b_prev_hash ->
        validate_chain([B2 | Rest], Valid)
    ;
        Valid = no
    ).
\end{lstlisting}

\subsection{Diagramme de la chaîne}

\begin{center}
\begin{tikzpicture}[
  block/.style={
    rectangle, rounded corners=5pt,
    minimum width=3cm, minimum height=2.5cm,
    draw=mercuryblue, fill=lightblue,
    line width=1.5pt, align=center, font=\small
  },
  arrow/.style={->, thick, mercuryblue}
]

% Bloc 0 - Genèse
\node[block, fill=lightgreen, draw=blockgreen] (g) at (0,0) {
  \textbf{Bloc \#0}\\
  \textit{Genesis}\\
  \footnotesize Nonce: 0\\
  \footnotesize Hash: 00ab...\\
  \footnotesize Prev: 0000...
};

% Bloc 1
\node[block] (b1) at (4.5,0) {
  \textbf{Bloc \#1}\\
  \textit{Alice→Bob}\\
  \footnotesize Nonce: 142\\
  \footnotesize Hash: 00f3...\\
  \footnotesize Prev: 00ab...
};

% Bloc 2
\node[block] (b2) at (9,0) {
  \textbf{Bloc \#2}\\
  \textit{Bob→Carol}\\
  \footnotesize Nonce: 387\\
  \footnotesize Hash: 00d7...\\
  \footnotesize Prev: 00f3...
};

% Bloc 3
\node[block] (b3) at (13.5,0) {
  \textbf{Bloc \#3}\\
  \textit{Carol→Dave}\\
  \footnotesize Nonce: 256\\
  \footnotesize Hash: 00e1...\\
  \footnotesize Prev: 00d7...
};

% Flèches avec label
\draw[arrow] (g.east) -- node[above,font=\tiny,mercurycyan]{prev\_hash} (b1.west);
\draw[arrow] (b1.east) -- node[above,font=\tiny,mercurycyan]{prev\_hash} (b2.west);
\draw[arrow] (b2.east) -- node[above,font=\tiny,mercurycyan]{prev\_hash} (b3.west);

% Difficulté
\node[below=0.5cm of b1, font=\footnotesize, color=blockorange]
  {PoW: Hash commence par "00"};

\end{tikzpicture}
\end{center}

\subsection{Résistance aux attaques}

\begin{center}
\begin{tabular}{p{4cm}p{5cm}p{5cm}}
\toprule
\textbf{Type d'attaque} & \textbf{Mécanisme de défense} & \textbf{Détection en Mercury} \\
\midrule
Falsification de données &
  Hash du bloc change → prev\_hash du suivant invalide &
  \code{validate\_chain} retourne \code{no} \\
\midrule
Double dépense &
  Consensus de la majorité du réseau &
  Chaîne la plus longue valide gagne \\
\midrule
Attaque 51\% &
  Coût computationnel prohibitif &
  Difficulté ajustée dynamiquement \\
\midrule
Replay attack &
  Chaque hash est unique (timestamp inclus) &
  Index + données + nonce = unicité \\
\bottomrule
\end{tabular}
\end{center}

%============================================================
\section{Comparaison des deux algorithmes}
%============================================================

\begin{center}
\begin{tabular}{lll}
\toprule
\textbf{Critère} & \textbf{Tri Bitonique} & \textbf{Minage Blockchain} \\
\midrule
Paradigme Mercury   & Récursion pure, pattern matching
                    & Récursion terminale, types algébriques \\
Complexité          & $\Oh{n\log^2 n}$ comparaisons
                    & $\Oh{2^d}$ hashs en moyenne \\
Déterminisme        & \code{det} (résultat unique)
                    & \code{det} (nonce trouvé déterministement) \\
Usage mémoire       & $\Oh{n\log n}$
                    & $\Oh{n}$ (liste de blocs) \\
Parallélisable      & Oui (réseau de comparateurs)
                    & Oui (mining pools) \\
Application réelle  & GPU sorting, bases de données
                    & Bitcoin, Ethereum, cryptomonnaies \\
\bottomrule
\end{tabular}
\end{center}

%============================================================
\section{Guide de compilation et d'exécution}
%============================================================

\subsection{Prérequis}

\begin{lstlisting}[language=bash, caption={Installation de Mercury (Ubuntu/Debian)}]
# Telecharger Mercury depuis mercury.cs.mu.oz.au
sudo apt-get install mercury-recommended

# Verifier l'installation
mmc --version
\end{lstlisting}

\subsection{Compilation}

\begin{lstlisting}[language=bash, caption={Compiler les deux programmes}]
# Tri bitonique
mmc --make bitonic_sort
./bitonic_sort

# Minage blockchain
mmc --make blockchain_mining
./blockchain_mining

# Avec optimisations
mmc -O2 --make bitonic_sort
mmc -O2 --make blockchain_mining
\end{lstlisting}

\subsection{Sortie attendue — Tri Bitonique}

\begin{tcolorbox}[colback=black!90,coltext=green!80,fontupper=\ttfamily\small]
=============================================\\
   TRI BITONIQUE --- Mercury Implementation\\
=============================================\\
--- Test 1 : Liste aleatoire ---\\
Entree  : [3, 7, 4, 8, 6, 2, 1, 5]\\
Sortie  : [1, 2, 3, 4, 5, 6, 7, 8]\\
Verification : OK\\
\\
--- Test 3 : Liste decroissante ---\\
Entree  : [16, 15, 14, 13, ..., 1]\\
Sortie  : [1, 2, 3, 4, ..., 16]\\
Verification : OK
\end{tcolorbox}

\subsection{Sortie attendue — Minage Blockchain}

\begin{tcolorbox}[colback=black!90,coltext=cyan!80,fontupper=\ttfamily\small]
[Minage bloc \#1] Donnees : "Alice envoie 50 BTM a Bob"\\
Recherche du nonce -> Trouve !\\
Nonce  : 142\\
Hash   : 00f3a8b2\\
--------------------------------------------------\\
Chaine valide -- Integrite confirmee\\
Simulation attaque -> Chaine INVALIDE -- Attaque detectee !
\end{tcolorbox}

%============================================================
\section{Conclusion}
%============================================================

Ces deux implémentations illustrent la puissance de \mercury{} pour des
algorithmes complexes :

\begin{itemize}
  \item \textbf{Tri Bitonique} : démontre la récursion pure, la manipulation
        de listes, le pattern matching et la déclaration de modes — principes
        fondamentaux de \mercury{}.

  \item \textbf{Minage Blockchain} : exploite les types algébriques, la
        récursion terminale (optimisée par le compilateur), et la séparation
        propre entre logique pure et entrées/sorties via \code{io::di, io::uo}.
\end{itemize}

\begin{resultbox}
\mercury{} se distingue par sa \textbf{correction à la compilation} : le
système de types et de modes garantit l'absence d'erreurs d'unification et
de déterminisme avant l'exécution, ce qui en fait un choix robuste pour des
algorithmes critiques comme le tri haute performance et la sécurité
cryptographique.
\end{resultbox}

\vfill
\begin{center}
  \small\color{gray}
  Rapport généré — INF322 Université de Yaoundé I $\cdot$
  \mercury{} 14.01+ $\cdot$
  \today
\end{center}

\end{document}
